\documentclass[10pt,a4paper]{amsart}
\usepackage[margin=19mm]{geometry}
\usepackage{amsmath,amssymb,mathtools}
\usepackage[scr=rsfs,cal=euler]{mathalfa}
\usepackage{xfrac}
\usepackage[unicode,hidelinks,bookmarks=false]{hyperref}
\newcommand{\ie}{i.e.\ }
\newcommand{\cf}{cf.\ }
\newcommand{\resp}{respectively\ }
\newcommand{\Subsection}{\subsection}
\providecommand{\nobreakdash}{\nobreak\hbox{-}}
\setlength{\emergencystretch}{2em}
\multlinegap0pt
\usepackage{amssymb}
\usepackage{bm}
\newtheorem{theorem}{Theorem}
\newtheorem{corollary}{Corollary}
\title{Sharp weighted norm estimates for the positive Bergman operator: the two-parameter case}
\author{Benoit F. Sehba}
\date{Source: arXiv:2608.01159v2 (CC BY 4.0)}
\begin{document}
\maketitle

\noindent\emph{Source and licence.} Sharp weighted norm estimates for the positive Bergman operator: the two-parameter case. Version arXiv:2608.01159v2 (CC BY 4.0), distributed by arXiv under the Creative Commons Attribution 4.0 International licence, http://creativecommons.org/licenses/by/4.0/, as recorded in the arXiv licence metadata for this exact version and shown on the abstract page https://arxiv.org/abs/2608.01159v2. Full licence notice: \texttt{licenses/arxiv-2608.01159-CC-BY-4.0.txt}.\par\smallskip

\noindent\textit{Editorial scope.} Counted unit: the article's corollary cor:weaksharpextra (source0.tex lines 218--224). The article's complete original proof of it is delivered with it (source0.tex lines 226--242, one \texttt{proof} environment of 17 lines). No mathematics is written, repaired or re-derived: the statement and the proof are the article's own lines, in the article's own notation, and no retained line is substituted or re-flowed.  Retained uncounted as the source-local context the proof needs: lines 198--204.  Nothing was omitted from any retained span, so the package carries no omission record.  No retained line contains a citation, so the wrapper carries no bibliography and no bibliography entry.  No retained line contains a figure environment, an image-inclusion command or drawing code, so no image is delivered and the package declares no asset dependency.  Editorial formatting only, in addition: an AMS \texttt{amsart} preamble that restates the article's own declarations and macros used by the retained lines, namely the environments \texttt{theorem}, \texttt{corollary} on separate counters, together with the standard packages those lines need.  The retained environments print in this document as Theorem 1, Corollary 1 in order of appearance, whereas the article prints the same items with the numbering of its own full document.  The complete native source of the article as distributed in the arXiv e-print for this exact version is bundled unchanged as \texttt{sources/206601/source0.tex}.
\begin{theorem}\label{thm:sharpextra}
Let $T$ be an operator defined on a suitable class (e.g. $C_c^{\infty}$ or $\cup_p L^p(\omega^p\,\mathrm{d}V_{\alpha})$, $\alpha>-1$). Suppose $p_0,q_0$ are exponents and $\beta_0$ is a real number, with $1\le p_0\le q_0<\infty$, $\frac{2+\beta_0}{q_0}= \frac{2+\alpha}{p_0}$, and
$$\|\omega Tf\|_{q_0,\beta_0}\le c[\omega]^{\theta}_{B_{p_0,q_0,\alpha}}\|\omega f\|_{p_0,\alpha}$$
for all $\omega\in B_{p_0,q_0,\alpha}$ and some $\theta>0$. Then
$$\|\omega Tf\|_{q,\beta}\le c[\omega]^{\theta\max\{1,\frac{q_0}{p'_0}\cdot\frac{p'}{q}\}}_{B_{p,q,\alpha}}\|\omega f\|_{p,\alpha}$$
for all $p,q,\beta$ satisfying $1<p\le q<\infty$, $\frac{1}{p} - \frac{1}{q} = \frac{1}{p_0} - \frac{1}{q_0}$, and $\frac{2+\beta}{q} = \frac{2+\alpha}{p}$, and for all $\omega\in B_{p,q,\alpha}$.
\end{theorem}

\begin{corollary}\label{cor:weaksharpextra}
Suppose that for some $1\le p_0\le q_0<\infty$ and $\frac{2+\beta_0}{q_0}= \frac{2+\alpha}{p_0}$, the operator $T$ satisfies the weak-type $(p_0,q_0)$ inequality
$$\| Tf\|_{L^{q_0,\infty}(\omega^{q_0}\,\mathrm{d}V_{\alpha})}\le c[\omega]_{B_{p_0,q_0,\alpha}}^{\theta}\|\omega f\|_{p_0,\alpha}$$
for every $\omega\in B_{p_0,q_0,\alpha}$ and some $\theta>0$. Then $T$ also satisfies the weak-type $(p,q)$ inequality
$$\| Tf\|_{L^{q,\infty}(\omega^{q}\,\mathrm{d}V_{\alpha})}\le c[\omega]_{B_{p,q,\alpha}}^{\theta\max\{1,\frac{q_0}{p'_0}\cdot\frac{p'}{q}\}}\|\omega f\|_{p,\alpha}$$
for all $1<p\le q<\infty$ satisfying $\frac{1}{p}-\frac{1}{q} = \frac{1}{p_0} - \frac{1}{q_0}$ and $\frac{2+\beta}{q} = \frac{2+\alpha}{p}$, and for all $\omega\in B_{p,q,\alpha}$.
\end{corollary}

\begin{proof}
Set $T_{\lambda}f = \lambda\chi_{\{|Tf|>\lambda\}}$. For fixed $\lambda>0$, applying Theorem \ref{thm:sharpextra} to $T_{\lambda}$, we find a constant $c$, independent of $\lambda$, such that
\begin{eqnarray*}
\|\omega T_{\lambda}f\|_{q_0,\alpha} & = & \lambda|\{z\in\mathbb{R}_+^2:|Tf(z)|>\lambda\}|_{\omega^{q_0},\alpha}^{\frac{1}{q_0}}\\
 & \le & \|Tf\|_{L^{q_0,\infty}(\omega^{q_0}\,\mathrm{d}V_{\alpha})}\\
 & \le & c[\omega]_{B_{p_0,q_0,\alpha}}^{\theta}\|\omega f\|_{p_0,\alpha}.
\end{eqnarray*}
Thus, if $\omega\in B_{p,q,\alpha}$ and $\frac{1}{p} - \frac{1}{q} = \frac{1}{p_0} - \frac{1}{q_0}$, Theorem \ref{thm:sharpextra} shows that
$$T_{\lambda}: L^p(\omega^p\,\mathrm{d}V_{\alpha})\longrightarrow L^q(\omega^q\,\mathrm{d}V_{\alpha})$$
is bounded, with
$$\| \omega T_{\lambda}f\|_{q,\alpha}\le c[\omega]_{B_{p,q,\alpha}}^{\theta\max\{1,\frac{q_0}{p'_0}\cdot\frac{p'}{q}\}}\|\omega f\|_{p,\alpha},$$
where $c$ does not depend on $\lambda$. Hence
\begin{eqnarray*}
\| Tf\|_{L^{q,\infty}(\omega^{q}\,\mathrm{d}V_{\alpha})} & = & \sup_{\lambda>0}\|\omega T_{\lambda}f\|_{q,\alpha}\\
&\le& c[\omega]_{B_{p,q,\alpha}}^{\theta\max\{1,\frac{q_0}{p'_0}\cdot\frac{p'}{q}\}}\|\omega f\|_{p,\alpha}.
\end{eqnarray*}
\end{proof}

\end{document}
